\chapter{Evidence, Tracing, and Resumption}
\label{ch:evidence-tracing-resumption}
\label{ch:evidence}

Gert captures a complete, durable record of every execution for audit, compliance, incident response, and operational replay. This section specifies the append-only JSONL trace format, evidence capture model, per-step snapshot contract, run resumption semantics, replay mode, and OpenTelemetry integration. Together these mechanisms ensure that every run produces a self-contained artifact that answers three questions: \textit{what happened}, \textit{what evidence was collected}, and \textit{how can execution be recovered or replicated}.

% ------------------------------------------------------------------ %
\section{Trace File Format}
\label{sec:trace-file-format}
% ------------------------------------------------------------------ %

\subsection{File Location and Naming}

Each run produces a single trace file located at:

\begin{verbatim}
.runbook/runs/<run-id>/trace.jsonl
\end{verbatim}

The \texttt{<run-id>} is a URL-safe string assigned at run creation time (e.g.\ \texttt{20260418T143022-a3f9c1}). All run artefacts are co-located under the run directory:

\begin{verbatim}
.runbook/runs/<run-id>/
  trace.jsonl          # append-only JSONL event log
  snapshots/           # per-step state checkpoints
    step-0000.json
    step-0001.json
    ...
  attachments/         # binary evidence referenced from trace events
    <sha256>.bin
    ...
  run.yaml             # completion manifest (written on terminal event)
  annotations.jsonl    # operator notes (append-only)
  lock                 # PID lock (removed on clean exit)
\end{verbatim}

\subsection{JSONL Envelope}

Every line in \texttt{trace.jsonl} is a self-contained JSON object with a common envelope. Consumers MUST tolerate unknown fields to allow additive schema evolution.

\begin{verbatim}
{
  "event_id":  <string>,    // UUID v4 for correlation
  "sequence":  <int64>,     // monotonically increasing within this run
  "kind":      <string>,    // event kind (see §12.1.3)
  "timestamp": <RFC3339>,   // UTC timestamp with sub-second precision
  "run_id":    <string>,    // run identifier
  "path":      <array>,     // tree-path to the active node (may be null)
  "payload":   <object>     // kind-specific payload (see below)
}
\end{verbatim}

The \texttt{sequence} field provides a total order over events in a single run. It is assigned by the \texttt{RunStore.WriteTrace} implementation and is not trusted for cross-run ordering. The \texttt{event\_id} field is a UUID v4 assigned at event emission time for correlation with external systems.

\subsection{Event Type Catalogue}

The following event kinds are normative in the schema. All user-supplied string fields (command arguments, outputs, prompt responses) pass through the redaction pipeline before being written. Enum-constrained declarations (tool action \texttt{args.<name>}/\texttt{outputs.<name>}, runbook \texttt{inputs.<name>}/\texttt{outputs.<name>}, design/gert/sections/03-schema-vnext.tex \S\ref{para:enum-constraint}) carry their member list once per run, in the \texttt{plan.validated} record described in design/gert/sections/03d-parse-time-enforcement.tex \S The \texttt{ValidatedPlan} Type Contract; no other event in this catalogue --- including \texttt{tool/invoked} and \texttt{tool/completed} below --- repeats it. A redacted declaration's member list is replaced with the literal string \texttt{"<redacted>"} plus a \texttt{member\_count} integer wherever it would otherwise be emitted (design/gert/sections/03d-parse-time-enforcement.tex \S The \texttt{ValidatedPlan} Type Contract).

\subsubsection{\texttt{run/started}}
\label{sec:run-started-event-client-field}

Emitted once, as the first event in every trace file.

\begin{verbatim}
"payload": {
  "runbook_id":    <string>,  // runbook name/path
  "user":          <string>,  // actor identity
  "mode":          <string>,  // "real" | "replay" | "dry-run"
  "inputs":        <object>,  // resolved input values (redacted)
  "client":        <string>,  // "cli" | "server" | "mobile-ios" 
                              // | "mobile-android"
  "parent_run_id": <string>   // omitted if not an invoke child
}
\end{verbatim}

\subsubsection{\texttt{step/started}}

Emitted immediately before executing a step.

\begin{verbatim}
"payload": {
  "step_id":   <string>,  // stable step identifier
  "step_type": <string>,  // "cli" | "manual" | "tool" | "invoke"
  "step_name": <string>   // human-readable name from runbook YAML
}
\end{verbatim}

\subsubsection{\texttt{step/completed}}

Emitted after a step exits cleanly.

\begin{minted}{go}
"payload": {
  "step_id":    <string>,
  "exit_status": <int>,      // exit code for cli steps; 0 for others
  "duration_ms": <int64>,    // wall-clock duration in milliseconds
  "captures":   <object>,    // captured output variables (redacted)
  "evidence":   [            // zero or more evidence items
    {
      "name":   <string>,
      "kind":   "text" | "checklist" | "attachment",
      "value":  <string>,    // for kind=text
      "items":  <object>,    // for kind=checklist
      "path":   <string>,    // relative path for kind=attachment
      "sha256": <string>,    // hex SHA256 for kind=attachment
      "size":   <int64>      // bytes for kind=attachment
    }
  ]
}
\end{minted}

\subsubsection{\texttt{step/failed}}

Emitted when a step exits with a non-zero status or returns a runtime error.

\begin{minted}{json}
"data": {
  "step_id":    <string>,
  "error_type": <string>,  // "exit_code" | "timeout" | "governance"
                           // | "tool_error" | "invoke_failed"
  "error":      <string>,  // human-readable message (redacted)
  "exit_status": <int>,    // for error_type=exit_code
  "duration_ms": <int64>
}
\end{minted}

\subsubsection{\texttt{step/skipped}}

Emitted when a step is not executed due to branch evaluation or iterate exhaustion.

\begin{minted}{json}
"data": {
  "step_id": <string>,
  "reason":  "branch_not_taken" | "iterate_exhausted" | "dry_run"
}
\end{minted}

\subsubsection{\texttt{step/retrying}}

Emitted before each retry attempt (requires the retry policy feature).

\begin{minted}{json}
"data": {
  "step_id":       <string>,
  "attempt":       <int>,     // attempt number (1 = first retry)
  "max_attempts":  <int>,
  "backoff_ms":    <int64>
}
\end{minted}

\subsubsection{\texttt{tool/invoked}}

Emitted when the tool runtime dispatches a tool call.

\begin{minted}{json}
"data": {
  "tool_name":      <string>,
  "transport":      "stdio" | "jsonrpc" | "mcp",
  "correlation_id": <string>,  // UUIDv4, links to tool/responded
  "inputs":         <object>   // sanitized inputs
}
\end{minted}

\subsubsection{\texttt{tool/responded}}

Emitted when the tool call returns.

\begin{minted}{json}
"data": {
  "correlation_id": <string>,
  "exit_code":      <int>,
  "duration_ms":    <int64>,
  "truncated":      <bool>     // true if output was capped
}
\end{minted}

\subsubsection{\texttt{manual/prompted}}

Emitted when a manual step displays its prompt to the operator.

\begin{minted}{json}
"data": {
  "step_id":     <string>,
  "prompt_text": <string>,   // the prompt shown (redacted)
  "sla_seconds": <int>       // 0 if no SLA configured
}
\end{minted}

\subsubsection{\texttt{manual/completed}}

Emitted when the operator submits a response to a manual step.

\begin{minted}{json}
"data": {
  "step_id":      <string>,
  "response":     <string>,   // operator's text response (redacted)
  "attestation":  <string>,   // "operator_confirmed" | "skipped"
  "approver":     <string>,   // identity of responding actor
  "duration_ms":  <int64>
}
\end{minted}

\subsubsection{\texttt{governance/blocked}}

Emitted when a governance rule blocks execution.

\begin{minted}{json}
"data": {
  "step_id":          <string>,
  "rule":             <string>,  // allowlist | denylist | env_block
                                 // | output_redact | approval_gate
  "action_attempted": <string>,  // description of blocked action
  "policy":           <string>   // policy name or path
}
\end{minted}

\subsubsection{\texttt{governance/approved}}

Emitted when an approval gate is cleared.

\begin{minted}{json}
"data": {
  "step_id":         <string>,
  "approver":        <string>,  // identity of approver
  "approval_method": <string>,  // "interactive" | "mcp" | "api"
  "duration_ms":     <int64>
}
\end{minted}

\subsubsection{\texttt{run/completed}}

Terminal event for a successful run.

\begin{minted}{json}
"data": {
  "final_status": "passed" | "failed" | "skipped",
  "duration_ms":  <int64>,
  "step_count":   <int>,
  "passed":       <int>,
  "failed":       <int>,
  "skipped":      <int>
}
\end{minted}

\subsubsection{\texttt{run/failed}}

Terminal event for a run that halted due to an unrecoverable error.

\begin{minted}{json}
"data": {
  "failing_step_id": <string>,
  "error":           <string>,
  "duration_ms":     <int64>
}
\end{minted}

\subsubsection{\texttt{checkpoint}}

Internal event written by \texttt{RunStore.SaveCheckpoint} to link a snapshot file to the trace sequence. Not user-facing but required for resumption.

\begin{minted}{json}
"data": {
  "snapshot_file": <string>  // e.g. "step-0003.json"
}
\end{minted}

\subsection{Crash Safety}

The trace file is opened with \texttt{O\_APPEND} and each event is written as a single \texttt{write(2)} call followed by \texttt{fsync}. On POSIX systems, a write of fewer than \texttt{PIPE\_BUF} bytes (4096 bytes on Linux) to an \texttt{O\_APPEND} file is atomic. Because the JSON encoding of a single trace event is always well under 4096 bytes for typical runbooks, partial-line corruption on crash is avoided in practice. For large payloads (e.g.\ tool output), the runtime truncates at a configurable limit (default 4 KB per field) before encoding.

The Go implementation:

\begin{minted}{go}
// WriteTrace appends one event atomically and fsyncs.
func (d *DirRunStore) WriteTrace(event DurableEvent) error {
    p := filepath.Join(d.baseDir, "trace.jsonl")
    f, err := os.OpenFile(p, os.O_CREATE|os.O_WRONLY|os.O_APPEND, 0644)
    if err != nil {
        return fmt.Errorf("open trace: %w", err)
    }
    defer f.Close()

    data, err := json.Marshal(event)
    if err != nil {
        return fmt.Errorf("marshal event: %w", err)
    }
    data = append(data, '\n') // JSONL delimiter
    if _, err := f.Write(data); err != nil {
        return fmt.Errorf("write trace: %w", err)
    }
    return f.Sync() // durability guarantee
}
\end{minted}

On recovery, any line that fails \texttt{json.Unmarshal} is silently skipped. The last clean \texttt{checkpoint} event determines the recoverable state.

\subsection{Tamper Evidence}

The trace format supports optional HMAC-SHA256 signing of each event. When a signing key is configured (via \texttt{GERT\_TRACE\_KEY} or the workspace config), the runtime appends a \texttt{"sig"} field to the envelope containing the HMAC-SHA256 of the canonical JSON of all other fields. Verifiers can replay the trace and recompute signatures. This does not prevent deletion of events (append-only at the filesystem level provides that property) but detects modification of existing events.

% ------------------------------------------------------------------ %
\section{Evidence Capture}
% ------------------------------------------------------------------ %

\subsection{What Constitutes Evidence}

In the schema, ``evidence'' is any artefact that attests to a step's outcome:

\begin{itemize}
  \item \textbf{Command output} — stdout/stderr of a \texttt{cli} step, captured and embedded in \texttt{step/completed}.
  \item \textbf{Manual attestations} — operator-provided text, checklist responses, and explicit confirmations from \texttt{manual} steps.
  \item \textbf{Tool responses} — structured JSON returned by a \texttt{tool} step, captured in \texttt{step/completed.captures}.
  \item \textbf{File attachments} — binary or text files explicitly attached by an operator or produced by a step, stored in the \texttt{attachments/} directory and referenced by SHA256 digest.
\end{itemize}

\subsection{SHA256 Hashing}

All file attachments are content-addressed by their SHA256 digest. The \texttt{evidence} package exposes:

\begin{minted}{go}
// HashFile returns hex SHA256 and file size.
func HashFile(path string) (sha256hex string, size int64, err error)

// NewAttachmentEvidence creates an EvidenceValue with hash and size.
func NewAttachmentEvidence(path string) (*EvidenceValue, error)
\end{minted}

The hash is stored in the \texttt{step/completed} event alongside the relative path under \texttt{attachments/}. Consumers can verify integrity at any time by re-hashing the referenced file.

Text evidence fields (command output, operator responses) are not individually hashed — the overall trace event integrity is provided by the optional HMAC-SHA256 event signature described in \S~12.1.5.

\subsection{File Attachments}

When a step produces a binary artefact, the runtime:

\begin{enumerate}
  \item Copies the file to \texttt{.runbook/runs/<run-id>/attachments/<sha256>.<ext>}.
  \item Emits an \texttt{EvidenceValue} of kind \texttt{attachment} in the \texttt{step/completed} event with fields \texttt{path}, \texttt{sha256}, and \texttt{size}.
\end{enumerate}

Attachments are deduplicated by content: if two steps produce identical file content, they share one attachment file. The trace events may reference the same \texttt{sha256} value independently.

\subsection{Evidence Retention}

Evidence is retained for the lifetime of the run directory. The workspace-level retention policy (configurable in \texttt{gert.yaml}) specifies:

\begin{itemize}
  \item \textbf{\texttt{retention.runs}} — number of completed runs to retain (default: 50).
  \item \textbf{\texttt{retention.max\_age\_days}} — maximum age in days (default: 90).
  \item \textbf{\texttt{retention.attachments}} — whether to include attachments in retention (default: true).
\end{itemize}

The \texttt{gert gc} command enforces retention policy by removing the oldest run directories, deleting their \texttt{attachments/} content first.

% ------------------------------------------------------------------ %
\section{Run State Snapshots}
% ------------------------------------------------------------------ %

\subsection{Checkpoint Model}

A checkpoint is a complete serialisation of \texttt{RunState} written to \texttt{snapshots/step-NNNN.json} after every successfully completed step. The snapshot captures the minimum state needed to resume execution from the next step:

\begin{minted}{go}
type RunState struct {
    RunID             string
    RunbookPath       string
    Mode              string            // "real" | "replay" | "dry-run"
    StartedAt         time.Time
    Actor             string
    CurrentStepIndex  int               // flat-mode index (legacy compat)
    Path              TreePath          // tree-mode cursor (authoritative)
    Vars              map[string]string // resolved input bindings
    Captures          map[string]string // accumulated step outputs
    History           []*StepResult     // completed step records
    InvokeStack       []InvokeFrame     // nested invoke context stack
    IterateState      map[string]*IterateProgress
    ParallelSlots     map[string][]SlotState
    LastCheckpointSeq int64             // trace seq at checkpoint time
}
\end{minted}

\subsection{Atomic Snapshot Write}

Snapshots are written atomically: the runtime marshals the state to a \texttt{.tmp} file, fsyncs it, then renames it to the final path. An incomplete snapshot (e.g.\ from a crash mid-write) retains the \texttt{.tmp} suffix and is skipped during recovery.

\begin{minted}{go}
func (d *DirRunStore) SaveCheckpoint(state *RunState) (string, error) {
    filename := fmt.Sprintf("step-%04d.json", len(state.History))
    finalPath := filepath.Join(d.baseDir, "snapshots", filename)
    tmpPath   := finalPath + ".tmp"

    data, _  := json.MarshalIndent(state, "", "  ")
    os.WriteFile(tmpPath, data, 0644)          // write to temp
    tf, _ := os.Open(tmpPath); tf.Sync(); tf.Close() // fsync temp
    os.Rename(tmpPath, finalPath)              // atomic rename
    // ... then write checkpoint trace event
}
\end{minted}

\subsection{Minimum Resumption State}

The state required to resume a run from step $N$ is exactly the snapshot at \texttt{step-NNNN.json}. It contains:

\begin{itemize}
  \item All resolved input values (\texttt{Vars}).
  \item Outputs captured by completed steps (\texttt{Captures}).
  \item The tree cursor (\texttt{Path}) identifying which node executes next.
  \item Branch decisions already taken (encoded in \texttt{History} step statuses).
  \item Iterate progress counters and parallel slot states.
  \item The invoke stack for nested runbooks.
\end{itemize}

The runbook YAML itself is not stored in the snapshot; it is re-read from \texttt{RunbookPath} at resume time. Operators must not modify the runbook between the original run and resumption.

\subsection{Package Map in the Run Manifest}
\label{sec:package-resumption-manifest}

Per Barbara's ruling
\texttt{.squad/decisions/archive/barbara-tool-packages-architecture-ruling.md}
(AR-TP-7 §7.4--7.5, ratified 2026-08-09), the completion/run manifest
(\texttt{run.yaml}, \S\ref{sec:trace-file-format}) additionally records the
frozen package map produced at plan time
(\texttt{design/gert/sections/06-tool-runtime.tex} \S Tool Packages,
\S Digest, Evidence, Trace, Resume):

\begin{minted}{yaml}
# run.yaml (excerpt)
packages:
  - name: acme.incident-tools
    version: "1.4.2"
    root: vendor/acme-incident-tools
    digest: "sha256:9f2e1a..."
catalogDigest: "sha256:7bd4c2..."
\end{minted}

This is covered by the same HMAC trace chaining as every other manifest
field (\texttt{design/gert/sections/08-security-and-trust.tex}
\S Trace File Integrity). Package roots and digests are therefore part of
the \textbf{Minimum Resumption State} in exactly the same sense as
\texttt{Vars} and \texttt{Captures} above: resuming a run means resuming
against the same tool universe it originally planned against, not merely
the same step cursor.

% ------------------------------------------------------------------ %
\section{Run Resumption}
% ------------------------------------------------------------------ %

\subsection{The \texttt{gert exec --resume} Contract}

\begin{minted}{bash}
gert exec --resume <run-id>
\end{minted}

The runtime executes the following recovery sequence:

\begin{enumerate}
  \item Locate \texttt{.runbook/runs/<run-id>/}.
  \item Attempt to acquire the \texttt{lock} file. Fail if a live process holds it.
  \item Call \texttt{RunStore.LoadLatestCheckpoint()} — scan \texttt{snapshots/} descending, skip \texttt{.tmp} files, return the newest valid JSON.
  \item Re-open \texttt{trace.jsonl} for append.
  \item Advance \texttt{CurrentStepIndex} by one (or advance the tree \texttt{Path} cursor past the last completed node).
  \item Re-build the governance engine from the runbook's policy block.
  \item Resume the execution loop from the restored \texttt{Path}.
\end{enumerate}

The resumption does \emph{not} re-execute any step recorded in \texttt{History}. It continues from the step \emph{after} the last checkpoint.

\subsection{Package Resumption Contract}
\label{sec:package-resumption-contract}

Immediately after step 6 (rebuilding the governance engine) and before step
7 (resuming the execution loop), \texttt{gert exec --resume} performs the
following package-integrity sequence (AR-TP-7 §7.5):

\begin{enumerate}
  \item Recompute every package digest and the catalog digest from the
        current filesystem, using the construction in
        \texttt{design/gert/sections/06-tool-runtime.tex} \S Digest,
        Evidence, Trace, Resume.
  \item Compare the recomputed digests against the values recorded in the
        run manifest's package map (\S\ref{sec:package-resumption-manifest}).
  \item \textbf{Mismatch is a hard refusal} --- \texttt{PKG-009} --- and the
        run \textbf{stays resumable}. It does \textbf{not} partially resume,
        and it does \textbf{not} silently re-plan. Resuming a run against
        changed tool definitions would invalidate the evidence chain, which
        is the entire point of the trace.
  \item \texttt{--allow-package-drift} overrides the refusal. When used, the
        runtime \textbf{MUST} emit a \texttt{governance/packageDriftAccepted}
        trace event recording both digests and the operator identity
        (\texttt{design/gert/sections/07-runtime-events.tex}
        \S Package and Substitution Events). Drift acceptance is an
        auditable act, not a silent bypass.
  \item Missing packages on resume are \texttt{PKG-001}, never a fallback to
        a different tier.
\end{enumerate}

\subsubsection{Replay mode is non-fatal}

Replay mode (\S\ref{sec:replay-mode}) does not invoke tools, so package
digests are \textbf{recorded and compared but non-fatal}: a mismatch emits a
\texttt{replay/packageDrift} warning event and replay proceeds. Replay
determinism is defined by the scenario file, not by the on-disk catalog.

\subsection{Resumability by Step Type}

\begin{itemize}
  \item \textbf{\texttt{cli} steps} — safe to resume. The step is re-executed from scratch. Runbook authors should ensure cli steps are idempotent (e.g.\ \texttt{kubectl apply} over \texttt{kubectl create}).
  \item \textbf{\texttt{manual} steps} — the operator is re-prompted. The prior response (if any) is lost. The prompt is re-displayed and the operator must re-attest.
  \item \textbf{\texttt{tool} steps} — resumable if idempotent. If the tool call was in-flight when the process crashed (i.e.\ a \texttt{tool/invoked} event exists with no matching \texttt{tool/responded}), the runtime re-issues the call. Authors must declare \texttt{idempotent: true} on tools where safe; others are re-executed with a warning log.
  \item \textbf{\texttt{invoke} steps} — if the child run completed (its terminal event is in the parent trace), the invoke is treated as completed. If the child run was in progress, the child is itself resumed via the nested run ID stored in the \texttt{InvokeStack}.
  \item \textbf{\texttt{branch}/\texttt{iterate}} — structural nodes; not executed directly. Cursor advancement handles re-entry correctly.
\end{itemize}

\subsection{Idempotency Guidance}

Runbook authors should write resumable runbooks by following these guidelines:

\begin{itemize}
  \item Use \texttt{kubectl apply} and \texttt{terraform apply} instead of imperative create commands.
  \item Use \texttt{--if-not-exists} or \texttt{--create-or-replace} flags where available.
  \item For tool steps, declare \texttt{idempotent: true} in the tool definition when the operation is safe to call multiple times.
  \item For manual steps, write prompts as questions that remain valid on a second ask (``Has the database change completed?'' rather than ``Confirm you have started the change'').
  \item Avoid steps whose side effects depend on the current time or a random seed without an explicit override mechanism.
\end{itemize}

\subsection{Partial Step Resumption}

If a \texttt{tool/invoked} event has no matching \texttt{tool/responded} event, the tool call was interrupted mid-flight. On resumption:

\begin{enumerate}
  \item The runtime logs a warning identifying the orphaned correlation ID.
  \item If \texttt{idempotent: true}: the tool call is re-issued from the beginning.
  \item If \texttt{idempotent: false} (default): the step is treated as \textbf{failed} and execution follows the normal failure path (error branch, if defined, or halt).
\end{enumerate}

% ------------------------------------------------------------------ %
\section{Replay Mode}
\label{sec:replay-mode}
% ------------------------------------------------------------------ %

\subsection{Overview}

\begin{minted}{bash}
gert exec --mode replay --scenario <scenario.yaml>
\end{minted}

Replay mode executes a runbook against a set of pre-recorded command responses and evidence values. No real commands are executed and no real tools are invoked. The mode is used for:

\begin{itemize}
  \item \textbf{Testing} — validating that a runbook's branching logic, captures, and governance rules behave correctly for a given set of inputs and command outputs (see \S~13 Testing and Acceptance).
  \item \textbf{Demonstration} — showing a runbook's execution flow without real infrastructure.
  \item \textbf{Regression} — verifying that a refactored runbook produces the same trace as the original.
\end{itemize}

\subsection{Scenario File Format}

A scenario file is a YAML document that maps commands to recorded responses and step IDs to recorded evidence:

\begin{minted}{yaml}
commands:
  - argv: ["kubectl", "get", "pods", "-n", "prod"]
    stdout: "NAME          READY   STATUS    RESTARTS\npod-abc   1/1     Running   0\n"
    stderr: ""
    exit_code: 0

evidence:
  check-pod-health:          # step_id
    observation:             # evidence name
      kind: text
      value: "pod is healthy, proceeding"
\end{minted}

Commands are matched in order of declaration (first match wins). Unmatched commands fail the scenario unless \texttt{allow\_unmatched: true} is set.

\subsection{Replay Semantics}

In replay mode:

\begin{itemize}
  \item \texttt{cli} steps — the command executor returns the pre-recorded stdout/stderr/exit\_code.
  \item \texttt{manual} steps — the evidence collector returns the pre-recorded evidence values instead of prompting the operator.
  \item \texttt{tool} steps — tool invocations are intercepted; if a matching recorded response exists, it is returned. Otherwise the step fails.
  \item \texttt{invoke} steps — the child runbook is executed in replay mode using the same scenario (or a nested scenario file if \texttt{scenario\_file:} is specified on the invoke step).
  \item \texttt{tool} steps whose action is implemented by \emph{substitution} (\texttt{execute.kind: runbook}, design/gert/sections/06-tool-runtime.tex \S Action Substitution \S Replay semantics, \S\ref{sec:tool-substitution-replay}) are replayed like an \texttt{invoke} step, not like a process-backed \texttt{tool} step: the substitute runbook is executed in replay mode against the \emph{same} scenario, and each of its own \texttt{cli}/\texttt{tool} steps is matched individually against that scenario's \texttt{commands:} list. A single recorded response is not attached to the substituted action as a whole, because the action has no process-shaped invocation of its own to match.
  \item Governance rules — evaluated normally; replay does not bypass governance.
  \item \texttt{enum} validation (design/gert/sections/03-schema-vnext.tex \S\ref{para:enum-constraint}) — evaluated normally; replay does not bypass it. Any argument or output value answered from the scenario is a bound value and is checked against its declared \texttt{enum} identically to a live run; a scenario recording an off-enum value fails replay with the same \texttt{ENUM-008}/\texttt{ENUM-009} the live path would raise.
  \item Trace events — written normally; the replay trace can be compared against a golden trace for regression testing.
\end{itemize}

\subsection{Determinism Requirement}

For replay to produce the same outcome across runs, the runbook must be deterministic given fixed inputs and command outputs. Non-determinism sources to avoid:

\begin{itemize}
  \item Timestamps in expressions (use captured values instead of \texttt{time.Now()}).
  \item Random identifiers in step names or captures.
  \item File system state that is not reproduced by the scenario.
  \item External API calls in tool steps that are not replayed (tool steps should be declared \texttt{idempotent: true} and covered by scenario fixtures).
\end{itemize}

% ------------------------------------------------------------------ %
\section{OpenTelemetry Integration}
% ------------------------------------------------------------------ %

\subsection{Rationale}

The JSONL trace format is optimised for durability, human readability, and offline analysis. OpenTelemetry (OTel) spans are optimised for real-time observability, distributed tracing, and integration with platforms such as Jaeger, Zipkin, Datadog, and Grafana Tempo. Both serve complementary purposes and supports both simultaneously.

\subsection{Span Hierarchy}

When OTel is enabled, the runtime creates a span hierarchy that mirrors the runbook execution tree:

\begin{minted}{yaml}
Span: gert.run                    // root span, covers entire run
  Span: gert.step                 // one per step executed
    Span: gert.tool.invoke        // one per tool call within a step
  Span: gert.step                 // next step
  ...
\end{minted}

The Go implementation uses \texttt{go.opentelemetry.io/otel/trace}:

\begin{minted}{go}
func (e *Engine) runWithOTel(ctx context.Context) error {
    ctx, runSpan := otel.Tracer("gert").Start(ctx, "gert.run",
        trace.WithAttributes(
            attribute.String("gert.run_id",    e.State.RunID),
            attribute.String("gert.runbook",   e.Runbook.Meta.Name),
            attribute.String("gert.actor",     e.State.Actor),
            attribute.String("gert.mode",      e.State.Mode),
        ),
    )
    defer runSpan.End()
    return e.execute(ctx)
}

func (e *Engine) executeStep(ctx context.Context, step Step) error {
    ctx, stepSpan := otel.Tracer("gert").Start(ctx, "gert.step",
        trace.WithAttributes(
            attribute.String("gert.step_id",   step.ID),
            attribute.String("gert.step_type", step.Type),
            attribute.String("gert.step_name", step.Name),
        ),
    )
    defer stepSpan.End()
    // ...
}
\end{minted}

\subsection{Span Attributes}

The following attributes are set on all spans:

\begin{center}
\begin{tabular}{lll}
\hline
\textbf{Attribute} & \textbf{Spans} & \textbf{Description} \\
\hline
\texttt{gert.run\_id}     & all       & Run identifier \\
\texttt{gert.runbook}     & all       & Runbook name \\
\texttt{gert.actor}       & run       & Actor identity \\
\texttt{gert.mode}        & run       & Execution mode \\
\texttt{gert.step\_id}    & step, tool & Step identifier \\
\texttt{gert.step\_type}  & step      & Step type \\
\texttt{gert.tool\_name}  & tool      & Tool name \\
\texttt{gert.transport}   & tool      & Transport (stdio/jsonrpc/mcp) \\
\hline
\end{tabular}
\end{center}

Trace context is propagated into tool invocations via the W3C \texttt{traceparent} header where the transport supports it (HTTP-based MCP tools). For stdio tools, context is passed as an environment variable \texttt{OTEL\_TRACEPARENT}.

\subsection{Relationship to JSONL Trace}

OTel spans and the JSONL trace are complementary, not redundant:

\begin{itemize}
  \item The JSONL trace is the authoritative durable record. It contains evidence hashes, governance events, snapshot references, and redacted output that OTel spans do not.
  \item OTel spans are ephemeral, push-based, and optimised for latency-sensitive observability dashboards.
  \item The trace event \texttt{"seq"} is embedded as the \texttt{gert.trace\_seq} span attribute, enabling correlation between spans and trace events.
\end{itemize}

\subsection{Configuration}

OTel is opt-in and configured via standard OTEL environment variables:

\begin{verbatim}
OTEL_EXPORTER_OTLP_ENDPOINT=http://localhost:4317
OTEL_SERVICE_NAME=gert
OTEL_TRACES_SAMPLER=always_on
\end{verbatim}

When \texttt{OTEL\_EXPORTER\_OTLP\_ENDPOINT} is not set, the OTel SDK is initialised with a no-op tracer and zero overhead is incurred. There is no \texttt{--otel} flag; configuration is entirely via the environment. This ensures that CI and offline use work without modification.

% ------------------------------------------------------------------ %
\section{Implementation Notes for Go}
% ------------------------------------------------------------------ %

The key Go interfaces for this subsystem are:

\begin{minted}{go}
// TraceWriter is the interface the engine uses to record events.
// Implementations must be safe for concurrent use from a single goroutine
// (the engine's execution loop); multi-goroutine safety is optional.
type TraceWriter interface {
    Write(event DurableEvent) error
    Close() error
}

// RunStore manages all artefacts for a single run.
type RunStore interface {
    RunID() string
    WriteTrace(event DurableEvent) error
    ReadTrace() ([]DurableEvent, error)
    ReadTraceSince(afterSeq int64) ([]DurableEvent, error)
    SaveCheckpoint(state *RunState) (string, error)
    LoadLatestCheckpoint() (*RunState, int64, error)
    AcquireLock() error
    ReleaseLock() error
    WriteManifest(manifest *RunManifest) error
}

// EvidenceCollector is the interface for gathering evidence at manual steps.
type EvidenceCollector interface {
    Collect(ctx context.Context, stepID string,
        prompts []Prompt) (map[string]*EvidenceValue, error)
}
\end{minted}

The \texttt{DirRunStore} is the production implementation. A \texttt{MemRunStore} (storing events in a slice) is provided for unit tests where filesystem I/O is undesirable. The \texttt{ReplayEvidenceCollector} returns pre-recorded evidence for scenario tests.

% ------------------------------------------------------------------ %
\section{Run Ingest API}
% ------------------------------------------------------------------ %

The run ingest API enables mobile and offline executors to upload completed runs to a central gert server for audit, compliance, and operational visibility. The API accepts runs in the standard trace JSONL format and reconstructs them server-side.

\subsection{Endpoint: \texttt{POST /api/v1/runs/ingest}}

Accepts a stream of JSONL trace events (identical format to \texttt{trace.jsonl}).

\paragraph{Request format:}

\begin{verbatim}
POST /api/v1/runs/ingest
Content-Type: application/x-ndjson
Authorization: Bearer <token>

{"event_id":"...","sequence":0,"kind":"run/started",...}
{"event_id":"...","sequence":1,"kind":"step/started",...}
{"event_id":"...","sequence":2,"kind":"step/completed",...}
...
\end{verbatim}

\paragraph{Response:}

\begin{minted}{json}
{
  "run_id": "20260418T143022-a3f9c1",
  "status": "ingested",
  "events_received": 47,
  "attachments_pending": ["a3f9c1...", "b2e4d5..."]
}
\end{minted}

The \texttt{attachments\_pending} field lists SHA256 digests of attachments referenced in the trace but not yet uploaded.

\subsection{Endpoint: \texttt{POST /api/v1/runs/\{run-id\}/attachments/\{sha256\}}}

Uploads a single attachment file.

\paragraph{Request format:}

\begin{verbatim}
POST /api/v1/runs/20260418T143022-a3f9c1/attachments/a3f9c1...
Content-Type: application/octet-stream
Authorization: Bearer <token>

<binary file content>
\end{verbatim}

\paragraph{Response:}

\begin{minted}{json}
{
  "sha256": "a3f9c1...",
  "size": 204800,
  "status": "stored"
}
\end{minted}

The server verifies the SHA256 checksum of the uploaded content before storing. If the checksum does not match the URL path, the request is rejected with HTTP 400.

\subsection{Idempotency}

Both endpoints are idempotent. Uploading the same trace or attachment multiple times has no adverse effect. The server deduplicates attachments by SHA256 digest.

% ------------------------------------------------------------------ %
\section{Run Ingest API}
% ------------------------------------------------------------------ %

The run ingest API enables mobile and offline executors to upload completed runs to a central gert server for audit, compliance, and operational visibility. The API accepts runs in the standard trace JSONL format and reconstructs them server-side.

\subsection{Endpoint: \texttt{POST /api/v1/runs/ingest}}

Accepts a stream of JSONL trace events (identical format to \texttt{trace.jsonl}).

\paragraph{Request format:}

\begin{verbatim}
POST /api/v1/runs/ingest
Content-Type: application/x-ndjson
Authorization: Bearer <token>

{"event_id":"...","sequence":0,"kind":"run/started",...}
{"event_id":"...","sequence":1,"kind":"step/started",...}
{"event_id":"...","sequence":2,"kind":"step/completed",...}
...
\end{verbatim}

\paragraph{Response:}

\begin{minted}{json}
{
  "run_id": "20260418T143022-a3f9c1",
  "status": "ingested",
  "events_received": 47,
  "attachments_pending": ["a3f9c1...", "b2e4d5..."]
}
\end{minted}

The \texttt{attachments\_pending} field lists SHA256 digests of attachments referenced in the trace but not yet uploaded.

\subsection{Endpoint: \texttt{POST /api/v1/runs/\{run-id\}/attachments/\{sha256\}}}

Uploads a single attachment file.

\paragraph{Request format:}

\begin{verbatim}
POST /api/v1/runs/20260418T143022-a3f9c1/attachments/a3f9c1...
Content-Type: application/octet-stream
Authorization: Bearer <token>

<binary file content>
\end{verbatim}

\paragraph{Response:}

\begin{minted}{json}
{
  "sha256": "a3f9c1...",
  "size": 204800,
  "status": "stored"
}
\end{minted}

The server verifies the SHA256 checksum of the uploaded content before storing. If the checksum does not match the URL path, the request is rejected with HTTP 400.

\subsection{Idempotency}

Both endpoints are idempotent. Uploading the same trace or attachment multiple times has no adverse effect. The server deduplicates attachments by SHA256 digest.
